\begin{thebibliography}{34}
\bibitem{r1} Hinton, G., Vinyals, O. \& Dean, J. Distilling the knowledge in a neural network. arXiv 1503.02531 (2015).
\bibitem{r2} Cloud, A. et al. Language models transmit behavioural traits through hidden signals in data. \textit{Nature} \textbf{652}, 615--621 (2026).
\bibitem{r3} Betley, J. et al. Emergent misalignment: narrow finetuning can produce broadly misaligned LLMs. In \textit{Proc. 42nd International Conference on Machine Learning}, PMLR \textbf{267}, 4043--4068 (2025).
\bibitem{r4} Askin, B. et al. Emergent and subliminal misalignment through the lens of data-mediated transfer. arXiv 2605.12798 (2026).
\bibitem{r5} Schrodi, S., Kempf, E., Barez, F. \& Brox, T. Towards understanding subliminal learning: when and how hidden biases transfer. In \textit{International Conference on Learning Representations} (2026); arXiv 2509.23886.
\bibitem{r6} Blank, C., Bhatia, A., Rajamanoharan, S., Conmy, A. \& Nanda, N. Subliminal learning is steering vector distillation. arXiv 2606.00995 (2026).
\bibitem{r7} Zur, A. et al. It's owl in the numbers: token entanglement in subliminal learning. https://owls.baulab.info (2025).
\bibitem{r8} Brockers, V. C., Ventzke, R. D., Neuhaus, V., Hidalgo-Ogalde, B. \& Priesemann, V. Learning through noise: why subliminal learning works and when it fails. arXiv 2605.23645 (2026).
\bibitem{r9} Madl, T. Channel location constrains the auditability of subliminal learning. arXiv 2606.22019 (2026).
\bibitem{r10} Shah, A., Chooi, J., Ou, J. \& Feng, S. Covert influence between language models. arXiv 2606.04071 (2026).
\bibitem{r11} Aden-Ali, I., Golowich, N., Liu, A., Shetty, A., Moitra, A. \& Haghtalab, N. Subliminal effects in your data: a general mechanism via log-linearity. In \textit{Proc. 43rd International Conference on Machine Learning}, PMLR \textbf{306}, 517--543 (2026).
\bibitem{r12} Magistrali, I., Berdoz, F., Dauncey, S. \& Wattenhofer, R. Subliminal signals in preference labels. In \textit{Agents in the Wild Workshop at the International Conference on Learning Representations} (2026); arXiv 2603.01204.
\bibitem{r13} Yang, A. et al. Qwen3 technical report. arXiv 2505.09388 (2025).
\bibitem{r14} Grattafiori, A. et al. The Llama 3 herd of models. arXiv 2407.21783 (2024).
\bibitem{r15} Efron, B. Bootstrap methods: another look at the jackknife. \textit{Ann. Stat.} \textbf{7}, 1--26 (1979).
\bibitem{r16} Geva, M., Caciularu, A., Wang, K. \& Goldberg, Y. Transformer feed-forward layers build predictions by promoting concepts in the vocabulary space. In \textit{Proc. 2022 Conference on Empirical Methods in Natural Language Processing} 30--45 (2022).
\bibitem{r17} Zou, A. et al. Representation engineering: a top-down approach to AI transparency. arXiv 2310.01405 (2023).
\bibitem{r18} Zhang, F. \& Nanda, N. Towards best practices of activation patching in language models: metrics and methods. In \textit{International Conference on Learning Representations} (2024); arXiv 2309.16042.
\bibitem{r19} Meng, K., Bau, D., Andonian, A. \& Belinkov, Y. Locating and editing factual associations in GPT. In \textit{Advances in Neural Information Processing Systems} \textbf{35} (2022); arXiv 2202.05262.
\bibitem{r20} Wang, K., Variengien, A., Conmy, A., Shlegeris, B. \& Steinhardt, J. Interpretability in the wild: a circuit for indirect object identification in GPT-2 small. In \textit{International Conference on Learning Representations} (2023); arXiv 2211.00593.
\bibitem{r21} Wu, W., Wang, Y., Xiao, G., Peng, H. \& Fu, Y. Retrieval head mechanistically explains long-context factuality. In \textit{International Conference on Learning Representations} (2025); arXiv 2404.15574.
\bibitem{r22} Olsson, C. et al. In-context learning and induction heads. \textit{Transformer Circuits Thread} https://transformer-circuits.pub/2022/in-context-learning-and-induction-heads/ (2022).
\bibitem{r23} McDougall, C., Conmy, A., Rushing, C., McGrath, T. \& Nanda, N. Copy suppression: comprehensively understanding an attention head. arXiv 2310.04625 (2023).
\bibitem{r24} Elhage, N. et al. A mathematical framework for transformer circuits. \textit{Transformer Circuits Thread} https://transformer-circuits.pub/2021/framework/ (2021).
\bibitem{r25} Singh, A. K. \& Strouse, D. J. Tokenization counts: the impact of tokenization on arithmetic in frontier LLMs. arXiv 2402.14903 (2024).
\bibitem{r26} Minder, J. et al. Narrow finetuning leaves clearly readable traces in activation differences. In \textit{International Conference on Learning Representations} (2026); arXiv 2510.13900.
\bibitem{r27} Geva, M., Schuster, R., Berant, J. \& Levy, O. Transformer feed-forward layers are key-value memories. In \textit{Proc. 2021 Conference on Empirical Methods in Natural Language Processing} 5484--5495 (2021).
\bibitem{r28} Sun, M., Chen, X., Kolter, J. Z. \& Liu, Z. Massive activations in large language models. In \textit{First Conference on Language Modeling} (2024); arXiv 2402.17762.
\bibitem{r29} Chen, R., Arditi, A., Sleight, H., Evans, O. \& Lindsey, J. Persona vectors: monitoring and controlling character traits in language models. arXiv 2507.21509 (2025).
\bibitem{r30} Fisher, R. A. \textit{Statistical Methods for Research Workers} (Oliver and Boyd, 1925).
\bibitem{r31} Vaswani, A. et al. Attention is all you need. In \textit{Advances in Neural Information Processing Systems} \textbf{30} (2017).
\bibitem{r32} nostalgebraist. Interpreting GPT: the logit lens. \textit{LessWrong} https://www.lesswrong.com/posts/AcKRB8wDpdaN6v6ru (2020).
\bibitem{r33} Park, K., Choe, Y. J. \& Veitch, V. The linear representation hypothesis and the geometry of large language models. In \textit{Proc. 41st International Conference on Machine Learning}, PMLR \textbf{235}, 39643--39666 (2024).
\bibitem{r34} Rimsky, N. et al. Steering Llama 2 via contrastive activation addition. In \textit{Proc. 62nd Annual Meeting of the Association for Computational Linguistics} 15504--15522 (2024).
\end{thebibliography}
